\documentclass[../main.tex]{subfiles}
\begin{document}
%==========================================TIÊU ĐỀ TẬP========================================
\begin{table}[h]
  \begin{adjustwidth}{-1cm}{}
    \begin{tabular}{>{\centering\arraybackslash}p{6cm}|>{\raggedright\arraybackslash}p{12.5cm}}
      \multirow{5}{6cm}
      % Cột bên trái: ÉPISODE và số tập
      {\centering
        {\\[-40pt]
          \flaregothic\fontsize{20pt}{20pt}\selectfont
          \color{eptype}ÉPISODE\color{black}\\
          \burbank\fontsize{72pt}{72pt}\selectfont
          \color{epnum}86
          \\[8pt]
          \hrule
          \vspace{8pt}
          \flaregothic\fontsize{20pt}{20pt}\selectfont
          \color{eptype}SPÉCIAL\color{black}\\
          \burbank\fontsize{72pt}{72pt}\selectfont
          \color{epnum}XIV
          \\[18pt]
          % \flaregothic\fontsize{14pt}{14pt}\selectfont
          % \color{epattr}BIENVENUE, \uline{\color{Banpire} Mel} !
          \flaregothic\fontsize{18pt}{18pt}\selectfont
          \color{holofes}L'HOLOFÊTE ! \\ \color{epattr} {\sen Mini-histoire}\\
          \color{black}
        }
      }
       &
      % Dòng thứ nhất: tên tập tiếng Nhật
      {\fotskip\fontsize{18pt}{18pt}\selectfont これがホロライブの\ruby{年}{ねん}\ruby{末}{まつ}！【\ruby{短}{たん}\ruby{編}{ぺん}\ruby{集}{しゅう}】
      }   \\[6pt]
      % Dòng thứ hai: tên tập tiếng Anh
       & {\cafeteria\fontsize{18pt}{18pt}\selectfont The 2\textsuperscript{nd} Fes: Short Stories}                         \\[4pt]
      % Dòng thứ ba: tên tập tiếng Việt
       & {\vnmsans\fontsize{18pt}{18pt}\selectfont Gặp nhau cuối năm 2020: Chuyện chưa kể trong 2\textsuperscript{nd} Fes} \\[8pt]
      % Dòng thứ tư: Nhân vật xuất hiện
       & {\brian{\fontsize{14pt}{14pt}\selectfont Track used: Modern Day - Final Wave}
         }
      \\[8pt]
      % Dòng cuối cùng: Thời gian diễn ra của tập (thường sẽ là ngày phát hành)
       & {\continuum\fontsize{16pt}{16pt}\selectfont Ngày phát hành: 27/12/2020}                                           \\[36pt]
    \end{tabular}
  \end{adjustwidth}
\end{table}
%=========================================TRUYỆN CHÍNH========================================
\color{black}

\begin{center}
  \uline{(Tập này được làm nhằm phục vụ cho hậu buổi concert của \textit{hololive 2nd fes. Beyond the Stage} vào ngày 21 và 22/12/2020.)}
\end{center}

\subsection*{\phantomsection\label{ep-S13.I}}
\addcontentsline{toc}{subsection}{{\sen{-Histoire 1-}} {\caviar Cuốn sách BL}}
\stpt{-Histoire 1-}\stcap{Cuốn sách BL}

``Nè, Rushia\~{} Đọc thử cái này đi.'' Cô nàng Thuyền trưởng chìa ra một cuốn sách với tựa đề đầy nghi vấn: ``\uline{Cuốn sách BL - Thịt xông khói x Rau diếp.}''

\img{7-4a1}

Nét mặt Senchou như một kẻ dâm đãng chính hiệu, cô sán lại gần Cô đồng và thì thầm với giọng điệu đầy hứng thú: ``Đây là một cuốn sách quý về chuyện tình lãng mạn giữa thịt xông khói và cây rau diếp đấy\~{}''

Rushia lập tức đổ mồ hôi hột, trán nhăn lại như không thể hiểu nổi: ``Ờm\dots\ bà đang nói cái quái gì vậy-desu?''

Nhưng Marine chẳng buồn giải thích. Cô nhẹ nhàng lướt tay trên bìa sách, đôi mắt mơ màng như thể đang chiêm ngưỡng một tuyệt tác nghệ thuật. Giọng cô tiếp tục nhỏ nhẹ, nhưng lại khiến người ta rùng mình: ``Bà có biết không? Đây chính là sự hòa quyện giữa hai hương vị tuyệt vời. Những lát thịt xông khói dai dai, hòa quyện cùng những lớp rau diếp giòn tan\dots\ Và chúng còn được 'tắm' trong lớp nước sốt vừng đầy quyến rũ nữa!''

Nàng Pháp sư nheo mắt đầy hoài nghi, cả người cứng đờ trước những lời lẽ mờ ám của Senchou: ``N-Này! Bỏ tôi ra!''

Nhưng Thuyền trưởng chẳng hề có ý định dừng lại. Cô uốn éo lắc lư cái hông, đôi mắt long lanh như thể đang truyền bá một tôn giáo mới: ``Bà biết không, Rushia? Những thớ thịt mềm mại của thịt xông khói\dots\ Những lá rau diếp mọng nước, nhai vui tai đến mức kích thích giác quan\dots\ Tất cả hòa quyện trên bàn ăn như một cặp tình nhân nồng cháy!''

\img{7-4a2}

Kuon, người đang ngồi cách đó không xa, không thể tập trung viết báo cáo được nữa. Cậu từ từ ngẩng mặt lên, nhìn hai thành viên thần tượng đang\dots\ dây dưa theo nghĩa đen. Trong đầu cậu lúc này chỉ có một câu duy nhất: ``\vie{\dots Cái lờ gì thế?}''

Rushia lúc này đã lắp bắp không nói nên lời: ``Ý-Ý của bà là gì chứ!?''

Bất ngờ, Marine liền đứng thẳng dậy, áp sát thẳng mặt vào cô, rồi hét lên: ``\textbf{Buổi \hll\ 2nd Fes - Beyond the Stage! Tài trợ bởi Bushiroad!}''

Cả Kuon và Rushia cùng đồng thanh đồng thanh khó hiểu: ``Bà/Em đang nói cái củ kẹc gì thế!?''

Cuối cùng, không ai hiểu vì sao nàng ``Hải tặc'' lại hét lên tên chương trình mà họ đang biểu diễn, có lẽ vì họ sợ khán giả\dots\ mất tập trung.

\subsection*{\phantomsection\label{ep-S13.II}}
\addcontentsline{toc}{subsection}{{\sen{-Histoire 2-}} {\caviar Cuộc phỏng vấn kỳ lạ}}
\stpt{-Histoire 2-}\stcap{Cuộc phỏng vấn kỳ lạ}

Tại một góc bàn, Choco và Kuon ngồi đối diện với Sora, sẵn sàng cho buổi phỏng vấn đặc biệt.

Ở giữa họ, Pekora nằm gục xuống bàn như một cái xác không hồn. Trên đầu cô nàng thỏ còn áp một cái đồng hồ báo thức, từng cơn co giật nhẹ chạy dọc sống lưng như thể tinh thần đã rời khỏi thể xác.

Choco hào hứng tuyên bố: ``Bọn em sẽ bắt đầu cuộc phỏng vấn nha!''

Nàng Thần tượng nhẹ nhàng gật đầu đồng ý, nhưng chưa kịp chuẩn bị tinh thần thì câu hỏi đầu tiên đã ập đến.

``Câu hỏi đầu tiên: Chị có thích Landolt C\footnote{Công cụ kiểm tra thị lực của một người dưới dạng một tấm bảng có các chữ C, mỗi chữ C sẽ có một vạch từ bốn hướng: lên, xuống, trái, phải. Kích cỡ chữ C sẽ giảm xuống từ phần trên của bảng xuống phía dưới.} không?''

\img{7-4b1}

Kuon liếc mắt sang Choco, vẻ mặt đầy nghi hoặc: ``Anh thấy câu hỏi này kỳ quặc thế\dots\ Mà thôi kệ, em thử trả lời anh xem nào?''

Nàng Thần tượng hơi giật mình trước câu hỏi khó hiểu, nhưng vẫn ấp úng đáp lại: ``C-Chị nghĩ là nó cũng hay ho đấy\dots\ Nhưng mà khó chọn quá\dots''

Choco lập tức tiếp tục: ``Vậy giữa Landolt C và Pekora-sama - người đã làm việc tăng ca suốt 1024 ngày, chị thích cái nào hơn?''

Vừa nói, cô vừa chỉ tay về phía nàng Thỏ.

Pekora trông chẳng khác gì một xác sống, toàn thân rũ rượi, con mắt thì trông như mất hết sức sống. Đồng hồ báo thức trên đầu cô vẫn kêu tí tách, còn cô thì co giật nhẹ, như thể đã thoát ly khỏi thực tại.

\img{7-4b2}

Cậu Tổng đổ mồ hôi hột, nhìn cảnh tượng trước mặt mà hỏi nhỏ: ``Em ấy trông có ổn không vậy\dots?''

Sora sững sờ nhìn người đồng nghiệp của mình, cảm giác như mình đang bị đặt vào một tình huống sinh tồn khắc nghiệt.

``Hả? Ờ thì\dots\ chị thích Pekora-chan hơn,'' cô tỏ vẻ thương hai cho nàng Thỏ xanh.

Choco lập tức nheo mắt đầy ẩn ý: ``Ồ?''

Kuon lập tức lắc đầu: ``Anh nghĩ đó không phải là lựa chọn phù hợp đâu. Cứ nhìn em ấy đi là biết. Anh cho em chọn lại đấy.''

Nàng Thần tượng câm nín nhìn nàng Thỏ co giật liên tục trên bàn. Sau một hồi suy nghĩ (và một chút sợ hãi), cô đành gãi đầu, cười trừ rồi ``quay xe'' một cách dứt khoát: ``Thôi thì\dots\ chị chọn Landolt C vậy!''

Ngay lập tức, cả nàng Y tá và cậu quản lý đồng thanh hô vang: ``Chị/Em được nhận!''

Bất ngờ thay, ngay khoảnh khắc đó, Pekora như bị đánh thức bởi một thế lực vô hình, cô bật dậy như lò xo rồi hét to: ``5:50 sáng rồi! 5:50 sáng rồi!'' Động tác của cô không khác gì như một chiếc đồng hồ báo thức đang vang lên inh ỏi.

\img{7-4b3}

Vừa hét, nàng Thỏ vừa tự đập vào đầu mình, đôi tai cô cuộn tròn lại làm hai tai đồng hồ, trông như thể đang băt chước một cái đồng hồ thật.

Sora tròn mắt kinh ngạc: ``Chị không biết là Pekora-chan lại là người thích dậy sớm đấy!''

Kuon thở dài, khoanh tay cười chát: ``Rồi em sẽ phải quen với môi trường ở đây thôi. Chào mừng em đến với công ty.''

Còn nàng Y tá thì vỗ tay không ngớt, mặt đầy vẻ thích thú như thể vừa chứng kiến một màn kịch đầy bất ngờ.

\subsection*{\phantomsection\label{ep-S13.III}}
\addcontentsline{toc}{subsection}{{\sen{-Histoire 3-}} {\caviar Xông hơi và\dots\ mù tạt?}}
\stpt{-Histoire 3-}\stcap{Xông hơi và\dots\ mù tạt?}

Tại một góc hành lang, Aqua đứng chờ sẵn, vừa thấy bóng Shion xuất hiện, cô nàng đã vẫy tay gọi: ``Shion-chan, ta đi xông hơi đi!''

Bước cạnh nàng Phù thủy là Kuon, người đang rảnh rang không có việc gì làm.

Nàng Phù thủy nghiêng đầu: ``Với Kuon-senpai sao?''

Cậu Tổng khoanh tay, nhíu mày đầy nghi hoặc: ``Sao anh có cảm giác không được chào đón ở đây nhỉ\dots?''

Shion lập tức xua tay: ``Em có nghĩ thế đâu!''

Rồi cô nàng quay sang cô bạn thân: ``Mà, đi xông hơi à? Đi luôn!''

Kuon nhìn cô với ánh mắt khó hiểu: ``Chấp nhận dễ dàng đến vậy sao\dots?''

\begin{center}
  \uline{Bên trong phòng xông hơi.}
\end{center}

Cả ba ngồi yên trong bầu không khí nóng bức, mồ hôi lấm tấm trên trán. Họ vẫn còn đang mặc quần áo trên người, không hề muốn cởi bỏ ra vì sợ đối phương nhìn thấy. Cặp đôi ``Hành - Tỏi'' ngồi ở một phía, Kuon ở phía đối diện.

\img{7-4c1}

Shion và Kuon đã bắt đầu thở dốc, mồ hôi chảy như thác nước.

Cậu lau trán, than thở: ``Nóng thật\dots''

Nàng Phù thủy cũng không khá hơn, cô cố cử động nhưng cảm giác như toàn thân mình đang tan chảy: ``Em đang rã mồ hôi như thác nước này\dots''

Kuon nhìn sang Aqua, thấy cô nàng vẫn ngồi yên bình, không hề có dấu hiệu mệt mỏi.

Cậu bĩu môi than thở: ``Ai bảo mặc quần áo vào làm gì\dots'' để rồi cậu sực nhớ ra cậu đang không ngồi một mình: ``Mà, đây là phòng ghép mà nhỉ?''

Nàng Hầu gái lười biếng trả lời: ``Hết phòng riêng rồi nên mới vậy chứ\dots''

Shion cười nhẹ, thả lỏng người hơn: ``Xông hơi đúng là phê thật, đúng không nhỉ, Aqua\dots''

Cô vừa quay sang bên cạnh, định trò chuyện với cô bạn thân của mình thì bất chợt nhìn thấy có hiện tượng lạ.

Trên cơ thể nàng Cô hầu, mồ hôi đang chảy xuống. Nhưng thay vì trong suốt, nó có một màu vàng kỳ quái.

\img{7-4c2}

Shion trợn tròn mắt: ``Cậu ta đang đổ mồ hôi trộm\dots\ à không, là thác Niagara mới đúng!''

Cậu Tổng nhìn một lúc rồi hốt hoảng: ``\dots Không, anh thấy em ấy đang đổ mù tạt ra đấy!''

Cả hai người lập tức lo lắng, cảm giác không khí trong phòng càng thêm ngột ngạt.

Nàng Phù thủy thở dài, lắc đầu bất lực: ``Trời ạ, Aqua-chan, tôi đã bảo bao nhiêu lần là đừng có giấu mù tạt trong người nữa mà.''

Kuon nhìn cô nàng hầu gái với ánh mắt khó hiểu: ``Để thế làm chua người chứ làm cái gì hả, Hành tím?''

Trong lúc cậu Tổng còn đang tìm cách xử lý chỗ mù tạt mà Aqua tiết ra, Shion chợt nảy ra một ý tưởng: ``A, nói đến hành tím thì\dots''

Cô nhếch môi cười tinh quái: ``Nếu ta kẹp cậu ta với bánh mì kẹp thịt xông khói và rau diếp thì sao nhỉ?''

Cô nàng lấy ra một chiếc bánh mì kẹp cất sẵn trong người. Cậu con trai lập tức quay ngoắt sang nhìn cô: ``Em kiếm cái đó ở đâu ra vậy?''

\img{7-4c3}

Rồi cậu nhăn mày, có một cảm giác quen thuộc dâng lên trong đầu: ``\textit{Khoan đã\dots\ thịt xông khói với rau diếp á? Sao nghe quen vậy?}''

Shion không trả lời, chỉ lẳng lặng lấy chiếc bánh kẹp ra, hứng mồ hôi mù tạt từ Aqua rồi\dots\ cắn một miếng to! Cô nàng lập tức reo lên: ``Ngon quá!''

Kuon sững người nhìn cô đang ăn miếng bánh ngon lành: ``\vie{Vãi.} Vậy cũng được nữa\dots?''

Nhưng rồi, thấy cô nàng đang hướng chiếc bánh kẹp về phía cậu, cậu cũng tỏ vẻ tò mò: ``Này, anh thử một miếng được không?''

\subsection*{\phantomsection\label{ep-S13.IV}}
\addcontentsline{toc}{subsection}{{\sen{-Histoire 4-}} {\caviar Nghệ thuật vì đam mê}}
\stpt{-Histoire 4-}\stcap{Nghệ thuật vì đam mê}

Một buổi chiều mát mẻ lộng gió, Miko đang đi dạo phố - một điều mà cô nàng rất hiếm khi làm, theo sau cô là Mel, Pekora và Kuon.

\img{7-4d1}

Nàng Ma cà rồng nghiêng đầu thắc mắc: ``Không biết tại sao tự dưng Miko-chan lại muốn ra ngoài nhỉ?''

Nàng Thỏ chống cằm suy nghĩ: ``Pekora cũng tò mò lắm ha, peko.''

Cậu Tổng nhún vai, vẻ mặt thản nhiên: ``Anh lại thấy bình thường mà. Giam mình mãi trong phòng cũng tù túng, thỉnh thoảng cũng nên ra ngoài hít thở tí chứ.''

Bất chợt, một cơn gió mạnh thổi qua, làm mọi người giật mình. Nhưng người ngạc nhiên nhất lại là Choco-sensei, vì cơn gió vô tình hất tung váy của cô lên.

``Kya\~{}!'' nàng Y tá thét lên, vội vàng giữ váy lại.

Pekora ôm đầu: ``Gió mạnh quá, peko!''

Kuon ngước lên bầu trời, khẽ nhíu mày: ``Kiểu này trời sắp mưa rồi đấy\dots''

Mel đưa tay vuốt tóc, lắc đầu cảm thán: ``Những cơn gió mạnh bất chợt thế này quả là bất tiện cho những người mặc váy nhỉ.''

Cậu con trai gật gù: ``Ừ. Trừ những ai đó thực sự biến thái\dots\ như em ấy chẳng hạn.''

Cậu hất cằm về phía Miko, người hiện đang đeo kính râm, trên tay cầm một chiếc máy ảnh, lia ống kính khắp nơi. Nhưng thay vì chụp cảnh vật, nàng Trông đền đang tập trung chụp vào những góc\dots\ rất khó diễn tả.

\gif{7-4d2}{1}{8}

Nàng Thỏ tròn mắt, vội kéo tay Mel: ``Hai người thấy chưa!? Miko-chi trong nháy mắt đã lôi cái máy ảnh ra và chụp rất nhiều pô từ hướng lên trên của váy Choco-sensei kìa, peko!''

Nàng Dơi há hốc miệng: ``Ể!? Thật á!?''

Cậu Tổng khoanh tay, gật gù đầy cảm khái: ``Nhìn cái cách lăn xả của em ấy là biết em ấy hy sinh vì nghệ thuật thế nào rồi.''

Mel trợn mắt, giận dữ lay vai cậu: ``Nghệ thuật cái gì chứ! Anh không định ngăn Miko-chan kia sao!?''

Kuon thở dài, giọng đầy bất lực: ``Mấy đứa biết não Miko-chan chỉ toàn có eroge với mới carh khiêu dâm thôi. Ngăn thế đếch nào được.''

``\textbf{Cái gì chứ!?}'' cả hai cô gái còn lại thảng thốt, không thể tin nổi với những gì họ nghe.

Pekora vỗ trán, giọng vừa ngán ngẩm vừa thảng thốt ngạc nhiên khi sự bất ngờ từ nàng Vu nữ chưa dừng lại ở việc chụp ảnh. ``Giờ chị ấy còn đang lập trình một tựa game dating sim (giả lập hẹn hò) của chính tay chị làm kìa, peko!!!''

\img{7-4d3}

Kuon liếc sang Miko, nhướng mày hỏi: ``Nhưng mà có biết viết code không đấy? Trình độ ê-lít như em ấy còn chưa dùng nổi VS nữa\dots''

Miko không trả lời, cô chỉ cười như một tên gàn dở. Trong đầu cô giờ đây chỉ có những cảnh nhạy cảm, và bản thân cô lại đang chìm đắm vào sở thích ``quái thai dị dạng'' này.

Cậu Tổng nhấp một ngụm nước, bình thản hỏi: ``Có muốn anh phụ phần lập trình không?''

Ngay lập tức, nàng Vu nữ gật đầu lia lịa. Không hiểu bằng cách nào, một chiếc ghế đột nhiên xuất hiện ngay bên cạnh cô. Cậu thản nhiên ngồi xuống, bắt đầu thảo luận với Miko về trò chơi.

``Này, này! Cho Pekora tham gia với!'' cô nàng Thỏ cũng không nhịn được, sà vào hai người họ đang bắt đầu nói chuyện rất rôm rả.

Mel nhìn cảnh tượng trước mắt, biểu cảm pha trộn giữa ngỡ ngàng, giận dữ, chán chường và\dots\ một chút kinh tởm.

\subsection*{\phantomsection\label{ep-S13.V}}
\addcontentsline{toc}{subsection}{{\sen{-Histoire 5-}} {\caviar Nghệ thuật điêu khắc vì tình yêu}}
\stpt{-Histoire 5-}\stcap{Nghệ thuật điêu khắc vì tình yêu}

Trong khu rừng vắng, Flare và Rushia đang đi dạo, vừa đi vừa trò chuyện.

Nàng Dị giới cười đùa với tai nạn do mình gặp phải vừa rồi: ``Nó dính đến mức mà tôi không thể thở được à!''

Nàng Pháp sư gật gù, tò mò hỏi: ``Thế à\dots\ vậy bà đã thoát ra bằng cách nào?''

Đằng sau họ, một bóng người đang ẩn nấp trong bụi cây - Noel, với khuôn mặt đầy ganh tị. Kuon đứng phía sau thì khoanh tay, cười khành khạch.

\img{7-4e1}

Nàng Hiệp sĩ nghiến răng, mắt long sòng sọc, tay run run nắm lấy thân cây gần đó: ``Flare\dots\ Tại sao bà có thể vui vẻ nói chuyện với người khác ngoài tôi chứ?''

Kuon ở phía sau liền hắng giọng, rồi đột nhiên cất tiếng hát từ bài của cặp đôi đã hát phiên bản thứ hai cho ``Rô-cô-na'': ``Bởi vì em ghen, ghen, ghen, ghen mà\~{}''

Câu hát của cậu như đổ thêm dầu vào lửa, khiến cô nàng càng tức tối hơn. Cô lôi ra một cây búa, mặt hằm hằm: ``Giờ thì tôi tức rồi đấy!''

Cậu Tổng giật mình, lùi lại một bước: ``Thôi, bỏ búa xuống. Kẻo lại làm hại người ta bây giờ!''

Noel lườm cậu, thở dài: ``Em có định đánh cậu ta đâu.''

Kuon ngơ ngác: ``Hả?''

Cô hếch cằm đầy kiêu hãnh, rồi giơ lên một cái đục: ``Em định chỉ khắc Flare làm từ gỗ cho bản thân em thôi.''

Cậu con trai đứng hình vài giây, sau đó phì cười: ``Tất nhiên là không thể rồi. Tính cách của em ấy hiền như hoa thế mà\dots''

Nàng Hiệp sĩ cúi xuống, lẩm bẩm với ánh mắt sáng rực: ``Heh\dots\ Rồi chúng ta sẽ cùng đi tận hưởng tuần trăng mật với nhau sau khi tôi khắc bà xong nha, Flare\~{}''

Nụ cười mất nhân tính của cô khiến cậu Tổng chỉ biết cười xòa, tự nhủ tốt nhất không nên xen vào nữa.

Nhưng ngay khi Noel định vung búa chặt cái cây đầu tiên để bắt đầu tác phẩm điêu khắc, một giọng nói quen thuộc vang lên từ phía xa: ``Kia chả phải Noel à! Bà đang làm gì đấy?''

Nàng Hiệp sĩ giật bắn người, búa suýt rơi khỏi tay. ``F-Flare!?''

\img{7-4e2}

Kuon khoanh tay, nửa đùa nửa thật: ``Cầu được ước thấy nha.''

Thấy thế, Noel lúng túng, ấp úng, mặt đỏ bừng vì không ngờ cô bạn thân của mình sẽ tới. Cô hoa tay múa chân, cố tìm một lời bao biện trong hành động dường như vô hại của mình:  ``K-Không phải như bà nghĩ đâu! Chỉ là\dots\ tôi muốn khắc họa thân bà thôi!''

\gif{7-4e3}{1}{13}

Cậu Tổng gật gù, nhún vai tiếp lời: ``Ít nhất là những gì em ấy muốn làm, Flare-chan.''

Nàng Dị giới nghiêng đầu khó hiểu: ``Ủa? Em sao? Ơ, thế là thế nào? Anh nói rõ hơn được không?''

\subsection*{\phantomsection\label{ep-S13.VI}}
\addcontentsline{toc}{subsection}{{\sen{-Histoire 6-}} {\caviar Liều thuốc mèo}}
\stpt{-Histoire 6-}\stcap{Liều thuốc mèo}

Tại văn phòng, Mio bất chợt ôm đầu, mặt xanh lè như tàu lá chuối, đôi mắt chỉ còn có những hình xoáy, như thể đang phát điên: ``\textbf{Á! Mình không có đủ mèo rồi! Mình cần một liều điều trị về mèo thôi!}''

Cậu Tổng đứng bên cạnh, mồ hôi túa ra: ``Anh chả hiểu ý em nói là gì, nhưng văn phòng mình làm gì có mèo đâu? Hay là bản năng Sói nguyên thủy của em lại lên cơn rồi?''

Vừa dứt lời, một giọng nói quen thuộc vang lên: ``Em có một bài dã thú đây!''

\img{7-4f1}

Okayu xuất hiện từ cửa, khuôn mặt vui thú nhưng ánh mắt lóe lên sự tinh quái. Cậu quản lý khẽ giật mình, rồi vỗ trán: ``À, anh quên mất là ta có một đứa Mèo ở đây.''

Cậu quay sang nàng Miêu, tò mò hỏi: ``Thế em định làm gì để `chữa bệnh' cho Sói đen-chan đây?''

Cô nàng không nói gì, chỉ lẳng lặng đẩy vào trong một Korone đang xoay vòng vòng như một chiếc ghế xoay, với cô nàng đang nằm ngửa, hai chân đá lên cao và hai tay co quặp lại.

\gif{7-4f2}{1}{27}

``\textbf{WowWowWowWowWowWowWowWow!}''

Kuon nhíu mày, chớp mắt vài lần trước cảnh tượng quái dị trước mặt: ``Thế quái nào em ấy có thể quay được mà không thấy bị chóng mặt nhỉ? Mà, đúng hơn là, em ấy xoay từ khi nào và bằng cách nào vậy?''

Okayu dường như bơ luôn câu hỏi cậu, mà chỉ đứng khoanh tay, nghiêm túc nói: ``Đừng có trốn tránh nữa, Koro-san!''

Korone tiếp tục xoay vòng vòng, hét lên: ``Tôi không thể dừng lại được!''

Cô cứ thế quay thêm mấy vòng nữa, nhưng rồi chậm dần, chậm dần\dots\ và cuối cùng dừng hẳn, không hề tỏ vẻ lảo đảo gì. Nàng Mèo tím thấy thế thì thở phào nhẹ nhõm lau mồ hôi: ``Thế là xong.''

Nhưng cậu con trai thì vẫn lắc đầu: ``Ờm, anh nghĩ em chưa xong đâu.''

Okayu nghiêng đầu, chưa kịp phản ứng thì bất ngờ, nàng Sói đen bỗng nhảy vào, hét lên: ``\textbf{Đây là cách điều trị mèo đó!!}'', rồi bắt đầu điên cuồng nhảy múa như thể vừa trúng bùa độc gì đó. Trông cô nàng chẳng khác gì một người vừa hút xì ke ma túy, hay là một người đang bị phát dại.

\img{7-4f3}

Cả cậu quản lý và nàng Miêu đều im lặng nhìn cảnh tượng hỗn loạn trước mặt.

Cuối cùng, Kuon chán nản thở dài: ``Anh nghĩ em ấy bị phê đá rồi\dots''

\subsection*{\phantomsection\label{ep-S13.VII}}
\addcontentsline{toc}{subsection}{{\sen{-Histoire 7-}} {\caviar Cánh tay máy thu hút}}
\stpt{-Histoire 7-}\stcap{Cánh tay máy thu hút}

Roboco đang hì hụi dưới sàn, tự mình điều chỉnh lại một bộ phận trong cánh tay. Ở phía sau, Aki đứng khoanh tay quan sát với ánh mắt tò mò, còn Kuon thì nhìn với vẻ thận trọng.

``Roboco-san, chị sửa xong cánh tay chưa thế?'' nàng Dị giới cúi xuống, nhìn cánh tay đang sửa mà tò mò hỏi.

``Rồi, chị đã chỉnh nó thành một cây sáo thu hút rồi,'' nàng Người máy đáp lại, hoàn thành những công đoạn cuối cùng.

\img{7-4g1}

``Lần này lại bày trò mèo gì nữa đây?'' cậu nhướng mày tỏ vẻ nghi hoặc. ``Anh hy vọng là em không bày trò rắc rối gì nữa\dots''

``Công nghệ của chị lúc nào cũng thú vị thật đấy! Nhưng nó hoạt động kiểu gì vậy?'' Aki hỏi.

Nàng Người máy cười hì hì rồi giơ cánh tay lên, bên trong cánh tay giờ có một cơ chế nhìn như một cây sáo tinh xảo: ``Chị có thể phát ra sóng thu hút từ đây. Chỉ cần thổi một tiếng, đối tượng sẽ bị mê hoặc ngay lập tức!''

Cậu con trai nhướng mày, nhìn cô với ánh mắt không mấy tin tưởng: ``Cái này có thực sự cần thiết không?''

``Tất nhiên rồi!'' cô hí hửng đáp. ``Thử nghĩ xem, nếu chị rơi vào tình huống nguy hiểm mà không thể di chuyển, chỉ cần thổi sáo, ai đó sẽ chạy đến giúp ngay lập tức!''

``Nghe cũng hợp lý đấy!'' nàng Tóc bay đáp lại một cách ngây thơ, choáng ngợp với cây sáo này.

``Nghe hợp lý ở đâu vậy, Aki-chan\dots'' cậu lắc đầu, tặc lưỡi. ``Thế em không tính đến chuyện nhỡ đâu kẻ được chạy đến lại là một tay biến thái thì sao?''

Roboco hí hửng chuẩn bị thử nghiệm, nhưng vì bản tính PON sẵn có, cô vô tình kích hoạt chế độ thu hút mà không để ý.

Vài giây sau, từ hành lang, hàng loạt người không mặt không mũi, toàn thân chỉ có một màu đen xì bất thình lình đổ xô vào căn phòng. Cả Aki và Kuon đều ngơ ngác khi chứng kiến cảnh tượng trước mắt.

\img{7-4g2}

``Á! Một nhóm học sinh không thể kiềm chế được ham muốn bản thân đang lao về phía cánh tay chị ấy kìa!\footnote{Nguyên văn: ​リビドーを抑えきれない学生が集まってきたよ！ (tạm dịch: Có nhiều học sinh không thể kiềm chế được ham muốn tình dục đang lao vào tìm cánh tay của Roboco kìa!)}'' nàng Dị giới thảng thốt. Càng lúc, lượng người như vậy càng đông, khiến văn phòng dần dần chật ních.

``\dots Và lại PON như mọi khi nhỉ, Đại PON-chan!?'' cậu Tổng cười khẩy, ngán ngẩm nhìn dòng người đen kịt đang bâu lấy cô nàng Người máy đang ngơ ngác kia.

Roboco hoảng hốt khi thấy mình bị vây quanh bởi một đám người có những hành động quái dị. Họ bắt đầu nhảy xung quanh cô một cách kỳ lạ, tạo thành một vòng tròn như đang thực hiện nghi thức gì đó.

``\textbf{Em không có ngốc mà!!}'' nàng Rô-bốt đáp lại trong những tiếng hò reo của đám người kia.

Trong lúc cả ba còn chưa kịp tìm cách xử lý, từ phía đám đông, một gương mặt quen thuộc bất ngờ xuất hiện. Một cô gái với nụ cười gian xảo đang chầm chậm tiến lại gần với đôi mắt sáng rực đầy hưng phấn.

\img{7-4g3}

``Matsuri-chan!?'' Cả ba người chết sững trong giây lát.

``Không thể nào\dots\ Cái này mà cũng tác dụng đến cả `quân' của mình ư!?'' Kuon thảng thốt, không thể nào tin nổi.

Nàng Dị giới bàng hoàng: ``Ai đó nhanh chóng giải tán nhóm người này đi với!''

``A-Aki-chan, Kuon-senpai, giúp chị với!!!'' Roboco vọng lại, vừa vật lộn với đám người nhưng không ra người ngợm, vừa la hét trong tuyệt vọng.

Matsuri tiến thêm một bước, nụ cười càng lúc càng ranh mãnh hơn. Đám đông sau lưng cô cũng đang chờ đợi\dots\ mà chắc chỉ có trời mới biết.

\subsection*{\phantomsection\label{ep-S13.VIII}}
\addcontentsline{toc}{subsection}{{\sen{-Histoire 8-}} {\caviar Tên gọi}}
\stpt{-Histoire 8-}\stcap{Tên gọi}

Suisei gọi Ayame, người vẫn đang chăm chú vào bài tập, không buồn ngẩng lên. Bên cạnh cô, cậu con trai đang ngồi khoanh tay, mắt lơ đãng nhìn về phía chân trời.

Nàng Ca sĩ cầm một cái vợt bắt cá, hào hứng hỏi: ``Ayame-chan, bọn chị đang định đi săn khủng long đấy. Muốn tham gia không?''

\img{7-4h1}

Cậu nhìn cái vợt mà Suisei đang cầm, nghiêng đầu tỏ vẻ khó tin: ``Ý em là đi tìm xương khủng long đúng không? Với cả cái vợt đó thì làm được trò trống gì?''

Chưa kịp để hai người phản ứng, đằng sau họ, một bộ ba nghịch ngợ - Matsuri, Haato, và Fubuki - đang nhảy nhót loạn cả lên, khiến cậu quản lý phải vội vàng đứng dậy can thiệp. ``\textbf{Này này, đừng có nhảy lên bàn! Matsuri, đừng có xách váy Haachama lên! Còn Fubuki, đừng có giả tiếng khủng long nữa!}''

Sau một hồi vật lộn với ba kẻ gây rối, cậu cuối cùng cũng tạm giữ chân được họ. Nhưng chỉ trong tích tắc, họ đã lại phấn khích hét lên: ``\textbf{KHỦNG LONG!}''

Và thế là, vòng lặp hỗn loạn lại tiếp tục, làm cậu Tổng vừa đuổi theo vừa cố gắng giữ họ trật tự.

Nhưng rồi, giữa mớ hỗn độn đó, Subaru bất ngờ chỉ tay về phía Ayame, người vẫn không hề rời mắt khỏi cuốn vở. ``Yên nào! Ayame-chan đang bận học để chuẩn bị cho kỳ thi lấy bằng Máy tính (cầm tay) điện tử đấy!''

\img{7-4h2}

``Toán bấm máy à? À, SAT chứ gì?'' cậu khẽ dừng lại, hỏi gọn.

Nghe vậy, nàng Oni cuối cùng cũng đặt bút xuống, nhướng mày nhìn sang nàng Vịt với vẻ mặt không mấy hài lòng. Cô rướn người lên bàn, hạ giọng như thể sắp sửa sửa lưng đối phương: ``Ý của em là bài kiểm tra lấy bằng Hướng dẫn Trà đạo Nhật Bản cơ!''

Kuon nhíu mày. ``Bài kiểm tra gì nghe lạ thế?''

Nhưng chưa kịp hiểu chuyện gì, Subaru lập tức bật lại: ``Hả? Bài kiểm tra Hướng dẫn viên du lịch Internet á?''

``Bài kiểm tra Chuyên gia phục vụ và thử nếm khăn tắm!'' nàng Oni đáp lại, mặt bừng bừng tức giận vì đối phương không hiểu ý.

Cậu Tổng đơ mất mấy giây với những tên gọi lạ hoắc này. ``\dots Cái củ quẹc gì vậy?'' Cậu sau đó thở dài, ngán ngẩm nói: ``Thế này thì mấy đứa nên ôn bài kiểm tra `Phương trình Ngữ pháp và Cú pháp câu văn tiếng Nhật' đi, vì anh còn chả hiểu mấy đứa đang nói gì cả!''

Cậu vừa dứt câu, ba kẻ phá hoại đằng sau lại tiếp tục màn đồng thanh quen thuộc: ``Hai, ba! \textbf{Khủng long!!}''

``Lại còn cả ba cái tên quỷ này nữa!?'' cậu quay ngoắt lại, mặt méo xệch. ``Làm ơn cư xử bình thường giúp anh cái!''

\img{7-4h3}

Suisei chỉ cười nhàn nhã với sự hỗn loạn này. ``À\dots\ Đây đích thực là \textit{hololive} mà chúng ta hay biết.'' Kuon đáp lại, trong lúc vẫn đang ``úp sọt'' ba cô gái xuống bãi cát: ``Dù anh ghét phải nói điều này, nhưng anh phải đồng ý\dots''

Cậu ôm trán, thở dài, thầm nghĩ: ``\textit{Mình đã làm gì để bị vây quanh bởi một đám `khủng long' dở hơi thế này chứ!?}''

\subsection*{\phantomsection\label{ep-S13.IX}}
\addcontentsline{toc}{subsection}{{\sen{-Histoire 9-}} {\caviar \textit{hologra} cần gì để bùng nổ?}}
\stpt{-Histoire 9-}\stcap{\textit{hologra} cần gì để bùng nổ?}

Trong căn phòng họp yên tĩnh, Kuon đang ngồi cùng với A-chan và Sora. Trên bàn, một loạt giấy tờ về kế hoạch sản xuất \textit{holo no Graffiti} (\textit{hologra}) được bày ra, nhưng không ai có vẻ hào hứng lắm.

``Có ai có ý kiến gì không?'' nàng Đeo kính hỏi.

Cậu Tổng giơ tay, và A-chan gật đầu gọi cậu đứng dậy phát biểu. ``Chị đã có kế hoạch gì cho \textit{hologra} sắp tới chưa?''

Đồng nghiệp của cậu chống cằm, nhìn xuống tập giấy trước mặt, rồi thở dài. ``Coi bộ cũng nhiều lắm, nhưng phần lớn toàn là xàm xí không à\dots\ Chẳng có gì đặc sắc cả.''

Cậu con trai nhướng mày. ``Thì đặc sản của nó chả phải là sự xàm xí sao?''

A-chan khẽ cười, gật gù: ``Biết là thế nhưng\dots\ nó vẫn thiếu cái gì ấy.''

Trong lúc hai người nhân viên đang vặn óc suy nghĩ, Sora lúc này giơ tay lên, rồi lên tiếng: ``Sự nhiệt tình của mọi người?''

Nàng Đeo kính lắc đầu: `Không, chị không nghĩ thế. Nhiệt tình thì đúng là có thừa, nhưng\dots''

Nàng Thần tượng nghiêng đầu suy nghĩ, rồi thử đoán tiếp: ``Vậy thì là sự nguyên bản?''

Cậu con trai chống tay lên bàn, nhún vai: ``Anh không nghĩ vậy đâu. Nếu không phải là Yamazu viết kịch bản thì cũng là các thành viên tự biên tự diễn mà. Hay là thiếu người?''

A-chan xua tay: ``Từ sau 2nd Fes thì chúng ta sẽ tuyển thêm thành viên mới thôi. Chắc không phải cái đó đâu.''

Cả ba người thở dài chán nản. Không khí phòng họp trở nên trầm lặng, ai cũng đang cố gắng tìm ra điều còn thiếu sót.

Cậu con trai lầm bầm trong miệng: ``\hl{Chả nhẽ nó lại thiếu sự nổi bật để cho người nước ngoài cảm thấy thích thú?}''

Đột nhiên, đôi mắt A-chan sáng bừng lên như vừa được khai sáng. ``Đúng rồi! Chính nó!''

``Hả?''

A-chan đập mạnh tay xuống bàn, đầy phấn khích. ``Từ trước đến nay ta chỉ dùng những trò đùa trong nước thôi, vậy tại sao ta không lấy những trò vui hay những tác phẩm từ nước ngoài nhỉ?''

Nàng Thần tượng cũng vỗ tay đồng tình với cô ấy: ``Cũng có thể lắm! Như thế thì sẽ càng có nhiều người nữa biết về chúng ta!''

Nàng Đeo kính nắm tay lại đầy quyết tâm. ``Quá tuyệt! Cảm ơn cậu nha, Hikawa-san!''

Hai người hào hứng bàn bạc về những nội dung quốc tế có thể đưa vào \textit{hologra}, bỏ lại cậu quản lý ngồi đó với vẻ mặt đầy sốc.

Nhìn hai cô nàng vui vẻ bàn tán với nhau, cậu vừa tỏ vẻ ngạc nhiên, vừa\dots\ hãnh diện khi cậu có thể giúp ích cho cả công ty này theo một cách nào đó. ``\textit{Mình nói vớ nói vẩn thế cũng được nữa\dots\ Hay thật!}"

\subsection*{\phantomsection\label{ep-S13.X}}
\addcontentsline{toc}{subsection}{{\sen{-Histoire 10-}} {\caviar Nhìn lại một năm xàm xí}}
\stpt{-Histoire 10-}\stcap{Nhìn lại một năm xàm xí}

Kuon đang ngồi trong phòng chiếu của trường, cắm tai nghe và dán mắt vào màn hình máy tính. Những thước phim cũ của \textit{hologra} lần lượt hiện lên, khiến cậu bật cười không ngừng.

Lúc này, cửa phòng mở ra. Hai người bạn cùng lớp là Shunyo và Tokue bước vào, thấy cậu cười tít mắt nên cậu bạn ``Sadboy'' liền hỏi: ``\vie{Đại tướng xem cái gì mà cười cành cạch thế?}''

Cậu con trai tháo tai nghe ra, quay sang đáp: ``\vie{À, mấy cái thước phim lộn xào mà tôi quay cùng với nhóm thần tượng holo ấy mà.}''

Cậu bạn Hiến máu nhướng mày ngạc nhiên, khoanh tay lại: ``\vie{Ý ông là cái nhóm như cái rạp xiếc ấy á?}''

Cậu Tổng gật đầu chắc nịch: ``\vie{Không họ thì ai.}''

Shunyo tò mò hỏi tiếp: ``\vie{Hỏi vui tí chứ trong số này ông thích tập nào nhất?}''

Kuon ngồi nghĩ một lúc lâu, rồi nở một nụ cười gian xảo, đáp gọn lỏn: ``\vie{Tập `mì trôi ống nước'.}''

Cậu Hiến màu nhăn mặt: ``\vie{Là cái bìu gì thế?}''

Cậu chàng hay cười khoanh tay, nheo mắt: ``\vie{Sao lại thích nó thế, Đại tướng?}''

Cậu Tổng bật cười, vỗ bàn: ``\vie{Đơn giản. Đồ ăn miễn phí. Đồng thời còn bắt nạt được cả Haachama nữa.}''

``\vie{Thế nó như nào?}'' cả hai cậu bạn cùng lớp đồng thanh. Kuon không trả lời mà chỉ hướng màn hình cho hai người cùng xem. ``\vie{Cứ xem đi thì biết.}''

\ct

Ban đầu, cả hai chỉ định xem qua một chút cho vui, nhưng chẳng biết thế nào mà bị cuốn theo cả loạt video. Những tập \textit{hologra} từ năm 2020 lần lượt trôi qua, hết tập này đến tập khác.

Thỉnh thoảng, Tokue lại cười sặc sụa, còn Shunyo thì hết lắc đầu lại bĩu môi khi họ nhìn lại chính mình (dù ít khi xuất hiện) trong đó.

``\vie{Hài thật đó!}'' Tokue vừa ôm bụng cười vừa nói.

``\vie{Đúng chứ? Mấy cái câu chuyện xàm xí này làm tôi dứt hết cả óc ra để giải quyết đấy.}''

Shunyo chậc lưỡi. ``\vie{Nhưng sao tôi thấy nó nhảm dị vậy\dots?}''

Kuon chỉ nhún vai, đáp gọn lỏn: ``\vie{Thì, \hll\ là vậy mà.}''

Cả ba cùng phá lên cười. Nhưng khi tiếng cười vừa lắng xuống, họ bỗng nhận ra xung quanh mình bỗng nhiên vắng tanh.

Tokue đảo mắt nhìn quanh phòng. ``Ơ\dots? Mọi người đi đâu rồi?''

Cậu Tổng quay ra nhìn cửa sổ. Bầu trời bên ngoài đã dần chuyển sang sắc tím nhạt, ánh đèn đường bắt đầu lóe sáng.

Cậu giật mình, đứng bật dậy. ``Vãi!? Đến giờ này rồi sao!?''

Shunyo cũng hốt hoảng: ``Thôi, anh em mình về nhanh kẻo trễ giờ xe!''

Không ai bảo ai, cả ba đứa con trai ba chân bốn cẳng thu dọn đồ đạc rồi cùng nhau chạy ào ra khỏi phòng, phóng như bay ra bến xe buýt.

Họ chuẩn bị trở về quê, đón một cái Tết Dương 2021 - năm Tân Sửu, khép lại một năm 2020 Canh Tý dương lịch đầy hỗn loạn nhưng cũng tràn ngập tiếng cười.

\end{document}